\section{Related work and the scope of the contribution}
\label{sec:related-work}

\paragraph{Adaptive mixtures of coins.}
Lee--Valiant is the closest direct antecedent \cite{lv}. In the extended
version, Sections~4--5 develop triangular single-coin estimators and
single-coin-to-population lower-bound decompositions. Section~7,
Theorems~4--5 and Lemma~22, optimizes an unbiased estimator when both
sign-conditional bias distributions, a nominal mixture weight, and a
per-coin cap are supplied. Its linear program optimizes information per
expected sample. Section~5.2 also reveals a coin's true bias after a cap
in a strengthened experiment.

Knowing our marginal magnitude law does not supply those conditional
distributions. For $0<\theta(q)<1$, the corresponding magnitude laws are
\[
 \mu_q^+(dd)=\frac{q(d)}{\theta(q)}\,\mu(dd),\qquad
 \mu_q^-(dd)=\frac{1-q(d)}{1-\theta(q)}\,\mu(dd).
\]
They change with the unknown function $q$. Our optimization must use one
observable procedure whose bias, moment, and expected work bounds hold
throughout $\mathcal Q_\rho$. The exact dual therefore retains separate
endpoint, moment-reference, and work-reference populations. The
additional statements concern that robust program, its comparison with
the full population experiment, and convergent upper and lower
certificates. Triangular flows, information/work optimization, adaptive
decomposition, and latent revelation are credited to their antecedents.

\paragraph{Gap entropy and sign testing.}
Chen--Li introduced entropy of normalized dyadic inverse-square gap
costs \cite[Definition~3.3]{cl16}. Chen--Li--Qiao proved the weighted
Sign-$\xi$ entropy inequality \cite[Lemma~4.1]{clq}; Appendix~D develops
its stopping-window argument, and Remark~D.4 explicitly covers Bernoulli
families away from their boundary. Our weighted sign-testing lemma is
an adaptation of this argument. In particular, the form
$A[1+\Ent(\pi)]$ and the statistical origin of its entropy factor are
established ideas.

The additional issue here is how much of a known magnitude law should
pay that entropy cost. The optimizing set in
\eqref{eq:scalar-selection} depends jointly on the law, rarity, and
budget. Its complement pays a truncated inverse-gap tail cost. The
upper construction implements this decomposition without observing
component membership; the lower construction respects both rarity
and mass saturation while combining disjoint components. These steps
are needed to obtain the all-law selection formula from the sign-testing
ingredient. Best-arm identification and latent population-proportion
estimation have different targets and sampling access.

\paragraph{Convex estimation, testing, and stopping.}
Juditsky--Nemirovski give constant-factor affine estimation under their
good-pair assumptions, with a Hellinger-affinity characterization
\cite[Theorem~3.1 and Proposition~3.1]{jn}. Polyanskiy--Wu characterize
quadratic minimax risk for an affine functional over a convex prior
class through a $\chi^2$ modulus for a fixed iid observation kernel
\cite[Theorem~1]{pw}. These results provide the relevant
estimation--testing duality. For a fixed capped one-unit policy, its
transcript is a fixed observation kernel. Optimizing that policy while
controlling expected work uniformly over $q$, and then comparing with
arbitrarily interleaved population algorithms, requires the additional
arguments given here. Strong convex duality and occupation-measure
formulations of optimal stopping are standard tools \cite{bv,cs}; the
four-prior formula is a specific application, not a new general duality
principle.

\paragraph{Baseline and narrower claims.}
The supplied manuscript and pinned version-5 development already
contain the balanced-law benchmark, rare-event phase results, the
variational $\Psi$ lower bound, and the critical quadratic-law entropy
correction \cite{original,repo}. The pinned critical proof also credits
Chen--Li--Qiao for its sign-entropy mechanism. Those results remain baseline contributions. The specialized ingredients
needed here are proved within this manuscript, with their attribution
preserved.

The principal added claim is the joint, arbitrary-law entropy--tail
characterization, together with a robust stopping-program formulation
that admits convergent certificates. The fixed-cap odd-moment
identification result is a secondary consequence of this model and
elementary moment interpolation; it is not presented as a general
mixture-identification principle. The targeted primary-source comparison
above identifies the closest checked statements and their hypotheses.
It does not certify publication priority.
